% !TEX root = ../../ctfp-print.tex

\lettrine[lhang=0.17]{如}{果我还没}能让你相信范畴论讲的全是态射, 那只能说明我没尽到职责. 既然下一个
主题是伴随 --- 它是由 hom-集之间的同构定义的 --- 回头梳理一下我们关于 hom-集这些积木的直觉是
合理的. 而且你会发现, 伴随为描述我们先前学过的许多构造提供了一套更通用的语言, 所以顺便复习它们
也有好处.

\section{函子}

首先, 你真的应该把函子想成态射的映射 --- Haskell 里 \code{Functor} 类型类的定义强调的就是这个
观点, 它整个围绕 \code{fmap} 展开. 当然, 函子也映射物件 --- 也就是态射的端点 --- 否则我们根本
没法谈保持组合. 物件告诉我们哪些态射对是可以组合的. 若两个态射要组合, 一个态射的目标必须等于
另一个态射的源. 所以, 如果我们想把态射的组合映到 \newterm{lifted 提升} 后的态射的组合, 那么
这些端点的映射基本上就已经确定了.

\section{交换图式}

态射的许多性质都是用交换图式来表达的. 如果某个态射可以用不止一种方式描述成其他态射的组合,
我们就得到了一个交换图式.

特别地, 交换图式构成了几乎所有普遍构造的基础 (初始物件和终止物件是显著的例外). 我们在积、余积、
各种别的 (余) 极限、指数物件、自由幺半群等的定义里都见过它.

积是一个普遍构造的简单例子. 我们挑两个物件 $a$ 和 $b$, 看看是否存在一个物件 $c$, 配上一对
态射 $p$ 和 $q$, 具有作为它们之积的普遍性质.

\begin{figure}[H]
  \centering
  \includegraphics[width=0.3\textwidth]{images/productranking.jpg}
\end{figure}

\noindent
积是极限的特例. 极限是用锥来定义的. 一个一般的锥由交换图式搭建而成. 这些图式的交换性可以换成
函子映射的恰当的自然性条件. 这样一来, 交换性就被降格成了自然变换这门高级语言的汇编语言.

\section{自然变换}

总的来说, 只要我们需要一个把态射映到交换方形的映射, 自然变换就非常方便. 自然性方形相对的两条
边, 是某个态射 $f$ 在两个函子 $F$ 和 $G$ 之下的映射; 另外两条边则是自然变换的分量 (它们也是
态射).

\begin{figure}[H]
  \centering
  \includegraphics[width=0.35\textwidth]{images/3_naturality.jpg}
\end{figure}

\noindent
自然性意味着, 当你走到``邻近''的分量上 (我说的邻近, 是指被某个态射连着), 你并不会与范畴或函子的
结构相对抗. 无论你是先用自然变换的分量跨过两个物件之间的间隙、再用函子跳到它的邻居, 还是反着来,
结果都一样. 这两个方向是正交的. 自然变换带你左右移动, 函子带你上下或前后移动 --- 可以这么说.
你可以把函子的\emph{像}想象成目标范畴里的一张薄片 (sheet), 而自然变换把对应于 $F$ 的这样一张
薄片映到对应于 $G$ 的另一张.

\begin{figure}[H]
  \centering
  \includegraphics[width=0.35\textwidth]{images/sheets.png}
\end{figure}

\noindent
我们在 Haskell 里见过这种正交性的例子. 在那里, 函子的作用修改容器的内容而不改变它的形状, 而自然
变换把原封不动的内容重新打包进另一个容器. 这两种操作的先后次序无关紧要.

我们见过极限定义里的锥被自然变换取代. 自然性保证了每个锥的那些边都交换. 尽管如此, 极限仍然是用
锥\emph{之间}的映射来定义的. 这些映射也必须满足交换性条件. (比如说, 积的定义里那些三角形必须
交换.)

这些条件同样可以由自然性来替代. 你可能记得, \emph{普遍}锥, 也就是极限, 被定义成这样两个函子
之间的自然变换 --- 一个是 (逆变的) hom-函子:
\[F \Colon c \to \cat{C}(c, \Lim[D])\]
另一个 (同样逆变的) 函子把 \emph{C} 中的物件映成锥, 而锥本身就是自然变换:
\[G \Colon c \to \cat{Nat}(\Delta_c, D)\]
这里 $\Delta_c$ 是常数函子, $D$ 是在 $\cat{C}$ 中定义图式的那个函子. 两个函子 $F$ 和 $G$
对 $\cat{C}$ 中态射的作用都有明确的定义. 恰好, $F$ 与 $G$ 之间的这个特定自然变换是一个
\emph{同构}.

\section{自然同构}

自然同构 --- 每个分量都可逆的自然变换 --- 是范畴论表达 ``两样东西相同'' 的方式. 这样一个变换的
分量必须是物件之间的同构 --- 一个拥有逆的态射. 如果你把函子的像想象成薄片, 那么自然同构就是这些
薄片之间一一对应且可逆的映射.

\section{hom-集}

那么, 态射又是什么呢? 相比物件, 它们确实有更多结构: 态射有两端, 物件没有. 但如果你固定源物件和目标
物件, 两者之间的态射就构成一个乏味的集合 (至少对局部小范畴来说如此). 我们可以给这个集合的元素
起 $f$ 或 $g$ 这样的名字, 以便把它们区分开 --- 但真正让它们彼此不同的, 到底是什么?

给定 hom-集中, 态射之间的本质区别在于它们与其他态射 (来自相接的 hom-集) 组合的方式. 如果存在
一个态射 $h$, 它与 $f$ 的组合 (无论是前组合还是后组合) 不同于它与 $g$ 的组合, 例如:
\[h \circ f \neq h \circ g\]
那我们就能直接 ``观察'' 到 $f$ 与 $g$ 的差别. 但即便这种差别无法直接观察到, 我们也可以借助函子
把 hom-集放大来看. 函子 $F$ 可能把这两个态射映成不同的态射:
\[F f \neq F g\]
在一个更丰富的范畴里, 相接的 hom-集提供了更高的分辨率, 比如:
\[h' \circ F f \neq h' \circ F g\]
其中 $h'$ 不在 $F$ 的像里.

\section{hom-集的同构}

许多范畴论构造都依赖 hom-集之间的同构. 但既然 hom-集只是集合, 它们之间一个普普通通的同构并说明
不了什么. 对有限集合, 同构只是说它们元素一样多. 若集合是无穷的, 那它们的基数必须相同. 但有意义的
hom-集同构必须考虑组合. 而组合牵涉的 hom-集不止一个. 我们需要定义横跨整簇 hom-集的同构, 还
需要施加一些与组合彼此相容的条件. 而一个 \newterm{natural 自然} 的同构, 恰好完全符合要求.

可 hom-集的自然同构又是什么? 自然性是函子之间映射的性质, 不是集合的性质. 所以我们真正在谈的,
是取值于 hom-集的函子之间的自然同构. 这些函子不只是取值于集合的函子, 它们对态射的作用是由恰当的
hom-函子诱导的. 态射被 hom-函子用前组合或后组合典范地映射 (取决于函子是协变的还是逆变的).

米田嵌入就是这样一个同构的例子. 它把 $\cat{C}$ 中的 hom-集映到函子范畴中的 hom-集; 而且是
自然的. 米田嵌入里的一个函子是 $\cat{C}$ 中的 hom-函子, 另一个把物件映到 hom-集之间的自然变换
构成的集合.

极限的定义同样是一个 hom-集之间的自然同构 (第二个照例在函子范畴里):
\[\cat{C}(c, \Lim[D]) \simeq \cat{Nat}(\Delta_c, D)\]
事实证明, 我们对指数物件的构造、对自由幺半群的构造, 也都能改写为 hom-集之间的自然同构.

这绝非巧合 --- 接下来我们会看到, 它们不过是伴随的不同例子, 而伴随正是被定义为 hom-集之间的
自然同构的.

\section{hom-集的不对称性}

还有一个观察会有助于我们理解伴随. hom-集一般是不对称的. hom-集 $\cat{C}(a, b)$ 往往与
hom-集 $\cat{C}(b, a)$ 大不相同. 这种不对称最彻底的例证, 就是把偏序当作范畴来看. 在偏序里,
从 $a$ 到 $b$ 的态射存在, 当且仅当 $a$ 小于或等于 $b$. 若 $a$ 与 $b$ 不同, 就不可能有反向
的、从 $b$ 到 $a$ 的态射. 所以, 如果 hom-集 $\cat{C}(a, b)$ 非空 (此时它意味着它是一个单例
集合), 那么除非 $a = b$, $\cat{C}(b, a)$ 必定是空的. 这个范畴里的箭头有明确的单向流动.

预序建立在未必反对称的关系上, 除了偶尔出现的环之外, 它也是 ``大体上'' 有方向的. 把任意一个范畴
想象成预序的推广, 在直觉上很方便.

预序是一个薄范畴 --- 所有 hom-集要么是单例, 要么是空的. 我们可以把一个一般的范畴想象成
``厚'' 一点的预序.

\section{挑战}

\begin{enumerate}
  \tightlist
  \item
        考虑自然性条件的几种退化情形, 画出恰当的图式. 比如说, 如果 $F$ 或 $G$ 中有一个函子把
        物件 $a$ 和 $b$ (也就是 $f \Colon a \to b$ 的两端) 都映到同一个物件, 例如
        $F a = F b$ 或 $G a = G b$, 会发生什么? (留意这样你就得到了一个锥或余锥.) 然后再考虑
        $F a = G a$ 或 $F b = G b$ 的情形. 最后, 如果你从一个绕回自身的态射出发 ---
        $f \Colon a \to a$ --- 又当如何?
\end{enumerate}
